\documentclass[11pt,letterpaper]{article}
\usepackage[T1]{fontenc}
\usepackage{lmodern,microtype}
\usepackage[margin=1in]{geometry}
\usepackage{amsmath,amssymb,amsthm}
\usepackage{xcolor}
\usepackage[colorlinks=true,linkcolor=blue!40!black,citecolor=blue!40!black,
            urlcolor=blue!40!black]{hyperref}
\hypersetup{
  pdftitle={The entropy photon-number inequality without an energy assumption},
  pdfauthor={Joseph M. Shunia},
  pdfsubject={Entropy of independent bosonic inputs under beam-splitter mixing},
  pdfkeywords={entropy photon-number inequality, bosonic Gaussian channels, von Neumann entropy}
}
\newtheorem{theorem}{Theorem}[section]
\newtheorem{lemma}[theorem]{Lemma}
\newtheorem{proposition}[theorem]{Proposition}
\newtheorem{corollary}[theorem]{Corollary}
\numberwithin{equation}{section}
\DeclareMathOperator{\Tr}{Tr}
\newcommand{\Fock}{\mathcal F_n}
\setlength{\emergencystretch}{1.5em}
\addtolength{\skip\footins}{6pt}

\title{The entropy photon-number inequality\\without an energy assumption}
\author{Joseph M. Shunia\thanks{Independent researcher, Ann Arbor, Michigan, United States.
  Email: \href{mailto:jshunia@gmail.com}{\texttt{jshunia@gmail.com}}.}}
\date{November 21, 2026}

\begin{document}
\maketitle

\begin{abstract}
We deduce the entropy photon-number inequality for two independent
bosonic inputs of finite von Neumann entropy from OpenAI's finite-energy
theorem. The result holds for any finite number of modes, even when
both inputs have infinite mean photon number. Applying the
same photon-number filter at both ports of the beam splitter gives an
operator bound on the retained marginal and hence an upper bound on its
entropy. At transmissivities strictly between zero and one, an input of
infinite entropy forces the output entropy to be infinite. We also
obtain the sharp minimum output entropy of a thermal attenuator at
prescribed input entropy, without an energy constraint.
\end{abstract}

\section{Introduction}

Mix two independent thermal beams of mean photon numbers $a$ and $b$
per mode. The retained beam is thermal with mean photon number
$\eta a+(1-\eta)b$ at transmissivity $\eta$. Guha, Erkmen, and
Shapiro conjectured that this pair minimizes the output entropy among
all independent inputs with the same respective entropies~\cite{ges}.
OpenAI proved the conjecture for any finite number of modes when both
inputs have finite mean photon number~\cite[Theorem~1.1]{oai}.
The energy hypothesis already excludes a simple one-mode state. In the
photon-number basis, give $k$ photons probability
\[
 p_k=\frac{6}{\pi^2(k+1)^2},\qquad k=0,1,\ldots.
\]
Its entropy is finite since $\sum_{k\ge1}\log(k+1)/(k+1)^2$ converges.
Its mean photon number is infinite since $k p_k$ is asymptotic to
$6/(\pi^2 k)$.

Write $S(\rho)=-\Tr(\rho\log\rho)\in[0,\infty]$ for the von Neumann
entropy, using natural logarithms and $0\log 0=0$. The entropy of a
one-mode thermal state of mean photon number $x$ is
\[
 g(x)=(1+x)\log(1+x)-x\log x,\qquad x\ge0.
\]
The function $g$ increases continuously from zero to infinity. Thus
$g^{-1}(S(\rho)/n)$ is the mean photon number per mode of the product
of $n$ identical thermal states having the same entropy as $\rho$.

\begin{theorem}\label{thm:main}
Let $A$ and $B$ be states of $n$ bosonic modes each, with $1\le n<\infty$
and $S(A),S(B)<\infty$. Mix the product state $A\otimes B$ at a beam
splitter of transmissivity $0\le\eta\le1$. If $a_j,b_j$ are the input
annihilation operators, let $C$ be the state of the modes
$c_j=\sqrt{\eta}\,a_j+\sqrt{1-\eta}\,b_j$, $1\le j\le n$. Then
\begin{equation}\label{eq:epni}
 S(C)\ \ge\ n\,g\!\left(
   \eta\,g^{-1}\!\left(\frac{S(A)}{n}\right)
   +(1-\eta)\,g^{-1}\!\left(\frac{S(B)}{n}\right)\right).
\end{equation}
\end{theorem}

The modes within either input may be entangled. An infinite value of
$S(C)$ satisfies \eqref{eq:epni} in the extended reals. These are the
hypotheses of Problem~418 in ProofAtlas~\cite{atlas}.
The Lean proof of Theorem~\ref{thm:main} is available at
\url{https://github.com/jshunia/autoresearch}.

The difficulty in deducing Theorem~\ref{thm:main} by approximation is
at the output port. If $C_r\to C$ in trace norm, lower semicontinuity
gives $S(C)\le\liminf S(C_r)$. To pass a lower bound on $S(C_r)$ to
$S(C)$ one needs an upper bound on the approximating entropies. We
obtain it by filtering both input ports by $r^N$, where $N$ counts
photons. Since the beam splitter conserves total photon number, the
same product filter acts on the two outputs. The filter on the
discarded port can then be removed at the cost of an operator
inequality on the retained port.

\section{An entropy estimate for filtered states}

For a positive trace-class operator $T$ of trace at most one, put
\[
 \Phi(T)=\Tr f(T),\qquad f(t)=t(1-\log t)\quad(0<t\le1),\qquad f(0)=0.
\]
The added trace term makes $f$ increasing on the whole state spectrum.
The following estimate is an elementary form of the spectral
domination used in Simon's convergence theorems for
entropy~\cite[Appendix, Theorems~A.2--A.4]{simon}.

\begin{lemma}\label{lem:entropy}
Let $\rho$ and $\sigma$ be states on a separable Hilbert space, let
$D$ be a contraction, and suppose $q\sigma\le D\rho D^*$ for some
$0<q\le1$. Then
\begin{equation}\label{eq:bound}
 q\bigl(1+S(\sigma)-\log q\bigr)\le 1+S(\rho).
\end{equation}
\end{lemma}

\begin{proof}
Write $\lambda_j(T)$ for the eigenvalues of a positive compact
operator, in decreasing order with multiplicity, padded by zeros in
the finite-rank case. The min--max principle gives
$\lambda_j(q\sigma)\le\lambda_j(D\rho D^*)$ for each $j$.
Taking $K=D\rho^{1/2}$, we have $KK^*=D\rho D^*$ and
$K^*K=\rho^{1/2}D^*D\rho^{1/2}\le\rho$. The nonzero eigenvalues of
$KK^*$ and $K^*K$ agree, so a second application of min--max gives
$\lambda_j(q\sigma)\le\lambda_j(\rho)$. Applying $f$ and summing the
nonnegative terms yields $\Phi(q\sigma)\le\Phi(\rho)$. Since
$\Phi(q\sigma)=q(1+S(\sigma)-\log q)$ and $\Phi(\rho)=1+S(\rho)$,
this proves \eqref{eq:bound}, including the case of an infinite trace
of $f(\rho)$.
\end{proof}

In particular, suppose $\sigma_r\to\rho$ in trace norm and
$q_r\sigma_r\le D_r\rho D_r^*$ with contractions $D_r$ and
$0<q_r\le1$ tending to one. If $S(\rho)<\infty$, \eqref{eq:bound} gives
$\limsup S(\sigma_r)\le S(\rho)$. Lower semicontinuity gives
$S(\rho)\le\liminf S(\sigma_r)$. Indeed, trace-norm convergence implies
convergence of each ordered eigenvalue by min--max, and Fatou's
lemma applies to the nonnegative terms
$-\lambda_j(\sigma_r)\log\lambda_j(\sigma_r)$. The same lower
bound suffices when $S(\rho)=\infty$. Thus in either case
\begin{equation}\label{eq:continuity}
 S(\sigma_r)\longrightarrow S(\rho)\quad\text{in }[0,\infty].
\end{equation}

\section{Photon-number damping at the two ports}

On the $n$-mode Fock space $\Fock$, let $P_k$ be the projection
onto the subspace of total photon number $k$, so
$N=\sum_{k\ge0}kP_k$. For $0<r<1$ and a state $\rho$ define
\begin{equation}\label{eq:filter}
 D_r=r^N,\qquad p_r(\rho)=\Tr(D_r\rho D_r),\qquad
 \rho_r=\frac{D_r\rho D_r}{p_r(\rho)}.
\end{equation}
Writing $w_k=\Tr(P_k\rho)$, we have
$p_r(\rho)=\sum_{k\ge0}r^{2k}w_k>0$ and $p_r(\rho)\uparrow1$ as
$r\uparrow1$. Each $\rho_r$ has finite mean photon number, since
\[
 \Tr(N\rho_r)=\frac{\sum_{k\ge0}k r^{2k}w_k}{p_r(\rho)}
 \le\frac{\sup_{k\ge0}k r^{2k}}{p_r(\rho)}<\infty.
\]
For trace-norm convergence, factor $\rho=\rho^{1/2}\rho^{1/2}$ and
use the Hilbert--Schmidt inequality to obtain
\[
 \|D_r\rho D_r-\rho\|_1\le2\|(I-D_r)\rho^{1/2}\|_2,\qquad
 \|(I-D_r)\rho^{1/2}\|_2^2=\sum_{k\ge0}(1-r^k)^2w_k\longrightarrow0.
\]
Hence $\rho_r\to\rho$ in trace norm. Taking $q_r=p_r(\rho)$ in
\eqref{eq:continuity} gives $S(\rho_r)\to S(\rho)$.

Let $U_\eta$ be the beam-splitter unitary on $\Fock\otimes\Fock$, let
$\omega=U_\eta(A\otimes B)U_\eta^*$, and take the first output port
as the retained one, so $C=\Tr_2\omega$. Write $C_r$ for the output
from the independently filtered inputs $A_r,B_r$, and set
$q_r=p_r(A)p_r(B)$. In general $C_r$ differs from the state obtained
by filtering $C$ alone.

\begin{proposition}\label{prop:output}
For every $0<r<1$,
\begin{equation}\label{eq:order}
 0\le q_r C_r\le D_r C D_r.
\end{equation}
Consequently $S(C_r)\to S(C)$ in $[0,\infty]$ as $r\uparrow1$.
\end{proposition}

\begin{proof}
The unitary $U_\eta$ preserves each eigenspace of
$N\otimes I+I\otimes N$. On the eigenspace of total number $m$,
$D_r\otimes D_r$ is multiplication by $r^m$. It therefore commutes
with $U_\eta$, and
\[
 q_r C_r=\Tr_2\bigl[(D_r\otimes D_r)\omega(D_r\otimes D_r)\bigr].
\]
Put $E_r=(I-D_r^2)^{1/2}$. For any bounded test operator $H$ on the
first port, the pairing of $D_r C D_r-q_r C_r$ with $H$ is
$\Tr[\omega(D_r H D_r\otimes E_r^2)]$. Thus
\[
 D_r C D_r-q_r C_r
 =\Tr_2\bigl[(D_r\otimes E_r)\omega(D_r\otimes E_r)\bigr]\ge0.
\]
By Lemma~\ref{lem:entropy}, \eqref{eq:order} gives the upper bound
$q_r(1+S(C_r)-\log q_r)\le1+S(C)$. The input products converge in
trace norm because $A_r\to A$ and $B_r\to B$ in that norm. Unitary
conjugation and partial trace give $C_r\to C$ in trace norm, so
\eqref{eq:continuity} applies with $q_r\to1$.
\end{proof}

\begin{proof}[Proof of Theorem~\ref{thm:main}]
Both filtered inputs have finite mean photon number. Their joint
output has finite total mean photon number by conservation, so each
marginal has finite entropy by the Gibbs bound~\cite[Lemma~2.2]{oai}.
The finite-energy inequality~\cite[Theorem~1.1]{oai} therefore gives
\begin{equation}\label{eq:regularized}
 S(C_r)\ge n\,g\!\left(
 \eta\,g^{-1}\!\left(\frac{S(A_r)}{n}\right)
 +(1-\eta)\,g^{-1}\!\left(\frac{S(B_r)}{n}\right)\right).
\end{equation}
By \eqref{eq:filter} and \eqref{eq:continuity} the input entropies
converge to $S(A)$ and $S(B)$. Letting $r\uparrow1$ and using the
continuity of $g$ and $g^{-1}$ on the right and
Proposition~\ref{prop:output} on the left proves \eqref{eq:epni}.
\end{proof}

\begin{corollary}\label{cor:infinite}
Suppose $A$ and $B$ are independent and $0<\eta<1$. If $S(A)=\infty$
or $S(B)=\infty$, then $S(C)=\infty$.
\end{corollary}

\begin{proof}
The filtered inputs have finite energy, so \eqref{eq:regularized}
applies for every $0<r<1$. If $S(A)=\infty$, then $S(A_r)\to\infty$.
Since $\eta>0$ and
$g^{-1}(s)\to\infty$ as $s\to\infty$, the right side of
\eqref{eq:regularized} tends to infinity regardless of $S(B_r)$.
Proposition~\ref{prop:output} gives the claim. If $S(B)=\infty$ use
$1-\eta>0$. At $\eta=0$ or $1$ the retained port carries just one
input, so the restriction on $\eta$ is necessary.
\end{proof}

\section{The minimum for a thermal attenuator}

Let $\tau_x$ denote the one-mode thermal state with mean photon number
$x\ge0$. Its photon counts have the geometric distribution of mean $x$,
with $\tau_0$ the vacuum. Let $\mathcal E_{\eta,b}$ be the one-mode
channel obtained by mixing the input with an independent copy of
$\tau_b$ and retaining the first output. For one mode, De Palma,
Trevisan, and Giovannetti proved the constrained minimum-output-entropy
theorem for phase-covariant Gaussian channels~\cite{dtg}. Multimode
results include the entanglement-breaking case~\cite{depalma} and
the heating range~\cite{bm}. OpenAI derived the following minimum over
finite-energy competitors~\cite[Corollary~1.2]{oai}.

\begin{corollary}\label{cor:thermal}
For $1\le n<\infty$, $a,b\ge0$, and $0\le\eta\le1$,
\[
 \min_{\substack{\rho\ \text{on }\Fock\\ S(\rho)=n g(a)}}
 S\!\left(\mathcal E_{\eta,b}^{\otimes n}(\rho)\right)
 = n\,g\bigl(\eta a+(1-\eta)b\bigr).
\]
The product thermal input $\tau_a^{\otimes n}$ attains the minimum.
\end{corollary}

\begin{proof}
The environment is $\tau_b^{\otimes n}$ and has entropy $n g(b)$.
Theorem~\ref{thm:main} gives the stated lower bound for every input
of entropy $n g(a)$. For attainment, the thermal mixing
identity~\cite[Proposition~2.3]{oai} gives
$\mathcal E_{\eta,b}(\tau_a)=\tau_{\eta a+(1-\eta)b}$, including
the vacuum cases. Taking the $n$-fold product gives equality.
\end{proof}


\section*{AI Assistance Disclosure}

OpenAI's GPT 6.2 Astra was used in an autoresearch loop and reported this discovery to the author.

\begingroup\small
\begin{thebibliography}{9}

\bibitem{ges}
S.~Guha, B.~I.~Erkmen, and J.~H.~Shapiro,
\emph{The entropy photon-number inequality and its consequences},
2008 Information Theory and Applications Workshop, IEEE (2008),
128--130. \href{https://doi.org/10.1109/ITA.2008.4601037}{doi:10.1109/ITA.2008.4601037}.

\bibitem{oai}
OpenAI, \emph{The entropy photon-number inequality}, OpenAI Math
Release preprint, September 24, 2026.
\href{https://github.com/openai/math/blob/4eaee87548250b2682affda7d05d947c05554bae/preprints/The-entropy-photon-number-inequality-September-24-2026/paper.pdf}{Manuscript}.

\bibitem{atlas}
ProofAtlas, \emph{Entropy photon-number inequality}, Problem~418,
2026. \href{https://www.proofatlas.ai/open-problems/\#problem-top500-omission-w045049-entropy-photon-number-inequality}{Problem statement}.

\bibitem{simon}
B.~Simon, \emph{Convergence theorems for entropy}, appendix to
E.~H.~Lieb and M.~B.~Ruskai, \emph{Proof of the strong subadditivity of
quantum-mechanical entropy}, Les rencontres physiciens-math\'ematiciens
de Strasbourg--RCP25 \textbf{19} (1973), expos\'e no.~5.
\href{https://www.numdam.org/item/RCP25_1973__19__A5_0/}{Numdam}.

\bibitem{dtg}
G.~De Palma, D.~Trevisan, and V.~Giovannetti,
\emph{Gaussian states minimize the output entropy of one-mode quantum
Gaussian channels}, Phys. Rev. Lett. \textbf{118} (2017), 160503.
\href{https://doi.org/10.1103/PhysRevLett.118.160503}{doi:10.1103/PhysRevLett.118.160503}.

\bibitem{depalma}
G.~De Palma, \emph{New lower bounds to the output entropy of multi-mode
quantum Gaussian channels}, IEEE Trans. Inform. Theory \textbf{65}
(2019), 5959--5968. \href{https://doi.org/10.1109/TIT.2019.2914434}{doi:10.1109/TIT.2019.2914434}.

\bibitem{bm}
S.~Beigi and H.~Mehrabi, \emph{From qubit tensorization to Gaussian
optimizers in bosonic systems}, preprint (2026).
\href{https://arxiv.org/abs/2610.28902}{arXiv:2610.28902}.

\end{thebibliography}
\endgroup

\end{document}
